\begin{proof}
\begin{enumerate}[label=\arabic*.]
\item Put $a=-256/3125$, and denote the whole reduced cusp fibre by
$W$.  The logarithmic sequence established in the proof of
\cref{night:P1A.2}, tensored with $K_M=h^*L^{-6}$ from
\cref{night:P1A.1}, is
\[
0\longrightarrow h^*L^5
\longrightarrow\Omega^1_M(\log W)\otimes K_M
\longrightarrow h^*(F\otimes L^{-6})\longrightarrow0.
\]
Here $K_{\widehat B}(a)=L^{11}$.  More precisely, in the meromorphic
$L$-frames of \cref{night:P1.15},
$\kappa=d\lambda/(\lambda-a)$ is $\lambda^9$ times the
$L^{11}$-frame: both sides have divisor $3D_0-D_\infty$.
The logarithmic sequence is exact in each chart
$q=z_0\cdots z_r$, where its first arrow sends $dq/q$ to
$\sum_{j=0}^r dz_j/z_j$; the quotient is the relative translation
cotangent bundle $h^*F$.  These arrows agree on the specified
monomial and deck transitions and descend through the free finite
root-chart actions.

\item For a summand $L^n$, the coarse degree is
$d(n)=\lfloor4n/5\rfloor+\lfloor-3n/4\rfloor$.
Projection formula and
$R^jh_*\mathcal O_M=\bigwedge^jG$ compute the two end terms.
For $h^*L^5$ the weights in relative degrees $j=0,1,2,3$ are
\[
\{5\},\quad\{5,4,3\},\quad\{4,3,2\},\quad\{2\}.
\]
Weights $4,5$ have degree zero and weights $2,3$ have degree $-1$.
Thus $h^q(h^*L^5)=(1,2,1,0,0)$.
The three weights of $F\otimes L^{-6}$ are $-4,-5,-6$.
After tensoring with $\bigwedge^jG$, every weight is between
$-9$ and $-4$.  We have $d(-9)=-2$ and $d(n)=-1$ for
$-8\le n\le-4$.  Weight $-9$ occurs once in each of relative
degrees $2$ and $3$, so
\[
h^q\bigl(h^*(F\otimes L^{-6})\bigr)=(0,0,0,1,1).
\]
The base has cohomological dimension one, so the two-column Leray
spectral sequences have no further differential.  Every connecting
map for the displayed logarithmic sequence has zero target.
Consequently
\[
h^q\bigl(\Omega^1_M(\log W)\otimes K_M\bigr)=(1,2,1,1,1).
\]

\item Let $\nu:Y\to W$ be the branch normalization.  The residue
sequence, with its actual branchwise maps, is
\[
0\longrightarrow\Omega^1_M\otimes K_M
\longrightarrow\Omega^1_M(\log W)\otimes K_M
\longrightarrow\nu_*\mathcal O_Y\otimes K_M|_W
\longrightarrow0.
\]
Since $K_M$ is pulled back from the base, its restriction to $W$ is
a constant line.  The complete smooth toric variety $Y$ has
$H^0(\mathcal O_Y)=\mathbb C$ and $H^{>0}(\mathcal O_Y)=0$
\cite[Theorem~9.2.3, p.~410, with divisor zero]{CoxLittleSchenck2011}.
The sole global section of the logarithmic middle term comes from
$h^*L^5$.  Its meromorphic expression is
$\lambda^{-5}\kappa$ tensored with the $L^{-6}$-frame: the
unique $L^5$ section is $\lambda^4$ times its frame, and
$\kappa$ contributes $\lambda^9$.  Its residue at $a$ is
$a^{-5}\ne0$.  Hence the map on global sections is an isomorphism
onto the residue line, and
\[
h^q(\Omega^1_M\otimes K_M)=(0,2,1,1,1).
\]
The pointwise perfect evaluation pairing
$T_M\otimes(\Omega^1_M\otimes K_M)\to K_M$ induces analytic
Serre duality
\cite[\S13, Theorem~4, p.~22; \S15, p.~23]{Serre1955Duality}.
It identifies $H^q(T_M)^\vee$ with
$H^{4-q}(\Omega^1_M\otimes K_M)$ and gives the stated tangent
cohomology dimensions.

\item We check the class of the actual parameter family on the
smooth fibre $A=h^{-1}(1)$.  In the fixed frame its period plane
has basis
\[
w_2,\quad w_1,\quad w_0+\frac{625\pi i}{96}\xi e.
\]
Differentiation in $\xi$, followed by projection to $G=H/F$,
gives the rank-one map $P:F\to G$ with sole nonzero matrix entry
$P_{13}=625\pi i/96$ in the ordered bases above and
$(e,w_{-1},w_{-2})$ of the quotient.
Differentiation in the base coordinate gives, by projecting the
connection in \cref{night:P1.15}, the matrix
\[
Q=\frac1{3125(1-a)}
\begin{pmatrix}
-78125/6+300r&-390625/48-375r&0\\
120&0&125/2\\
0&120&50
\end{pmatrix},
\qquad r=\frac{625\pi}{96}(-i\xi).
\]
Its determinant is $3/[625(1-a)^3]\ne0$.
The translation-valued \v Cech cocycles obtained by differentiating
the lattice gluing identify
$H^1(T_A)=\operatorname{Hom}(F,G)$; contraction with a relative
one-form sends the two cocycles to the displayed projected period
derivatives, in the same convention.

\item The normal tangent sequence on this fibre is
\[
0\longrightarrow T_A\longrightarrow T_M|_A
\longrightarrow\mathcal O_A\otimes T_{\widehat B,1}
\longrightarrow0.
\]
Its connecting map sends a base tangent vector to its
Kodaira--Spencer class $Q$.  As $H^0(\mathcal O_A)=\mathbb C$,
the kernel of $H^1(T_A)\to H^1(T_M|_A)$ is exactly
$\mathbb C Q$.  The nonzero matrix $P$ has determinant zero;
every nonzero scalar multiple of $Q$ has nonzero determinant.
Therefore the image of $P$ in $H^1(T_M|_A)$ is nonzero.
The global parameter Kodaira--Spencer cocycle restricts to this
image, so its class in $H^1(T_M)$ is nonzero.
The finite affine actions are fixed, and the cusp multipliers
$\exp(2\pi iC(q,\xi)\nu)$ vary holomorphically in $\xi$.
The same finite fan charts, with $C$ bounded on a small parameter
neighbourhood, give the proper local family of
\cref{night:P1.9,night:P1.10}.  Thus the cocycle just computed is
the derivative of this compact family.

\item Douady's versality and semiuniversality theorems
\cite[\S VII.8, pp.~598--599; \S VIII.1, pp.~599--600]{Douady1974}
give a miniversal germ with tangent space $H^1(T_M)=\mathbb C$.
Its local ring has a minimal presentation $\mathbb C\{u\}/I$,
where $I\subset(u)^2$.  The actual parameter disc classifies by
$u\mapsto g(s)$ with $g'(0)\ne0$, by the preceding step.
Substitution $\mathbb C\{u\}\to\mathbb C\{s\}$ is therefore an
isomorphism of convergent power-series rings.  Since this map factors
through $\mathbb C\{u\}/I$, it annihilates $I$; its injectivity
gives $I=0$.  The miniversal germ is reduced and smooth of dimension
one, and the classifying map is a local isomorphism.  This proves
miniversality of the given family at each member.
\end{enumerate}
\end{proof}
